\section{A representability criterion}
\label{section-representable}

\noindent
In this section we reformulate the glueing lemma
of Section \ref{section-glueing-schemes} in terms
of functors. We recall some of the material from
Categories, Section \ref{categories-section-opposite}.
Recall that given a scheme $X$ we
can define a functor
$$
h_X : \Sch^{opp}
\longrightarrow
\textit{Sets}, \quad
T \longmapsto \Mor(T, X).
$$
This is called the {\it functor of points of $X$}.

\medskip\noindent
Let $F$ be a contravariant functor from the category
of schemes to the category of sets. In a formula
$$
F : \Sch^{opp}
\longrightarrow
\textit{Sets}.
$$
We will use the same terminology as in Sites, Section
\ref{sites-section-presheaves}. Namely, given a scheme
$T$, an element $\xi \in F(T)$, and a morphism $f : T' \to T$
we will denote $f^*\xi$ the element $F(f)(\xi)$, and sometimes
we will even use the notation $\xi|_{T'}$

\begin{definition}
\label{definition-representable-functor}
(See Categories, Definition \ref{categories-definition-representable-functor}.)
Let $F$ be a contravariant functor from the category
of schemes to the category of sets (as above).
We say that $F$ is {\it representable by a scheme}
or {\it representable} if there exists a scheme $X$
such that $h_X \cong F$.
\end{definition}

\noindent
Suppose that $F$ is representable by the scheme $X$ and that
$s : h_X \to F$ is an isomorphism.
By Categories, Yoneda Lemma \ref{categories-lemma-yoneda}
the pair $(X, s : h_X \to F)$ is unique up to unique
isomorphism if it exists.
Moreover, the Yoneda lemma says that
given any contravariant functor $F$ as above
and any scheme $Y$, we have a bijection
$$
\Mor_{\text{Fun}(\Sch^{opp}, \textit{Sets})} (h_Y, F)
\longrightarrow
F(Y), \quad
s \longmapsto s(\text{id}_Y).
$$
Here is the reverse construction. Given any $\xi \in F(Y)$
the transformation of functors $s_\xi : h_Y \to F$
associates to any morphism $f : T \to Y$ the
element $f^*\xi \in F(T)$.

\medskip\noindent
In particular, in the case that $F$ is representable, there exists
a scheme $X$ and an element $\xi \in F(X)$ such that the corresponding
morphism $h_X \to F$ is an isomorphism.
In this case we also say {\it the pair $(X, \xi)$ represents $F$}.
The element $\xi \in F(X)$
is often called the {\it ``universal family''} for reasons that will become
more clear when we talk about algebraic stacks (insert future reference here).
For the moment we simply observe that the fact that if the pair $(X, \xi)$
represents $F$, then every element $\xi' \in F(T)$ for any $T$ is of
the form $\xi' = f^*\xi$ for a unique morphism $f : T \to X$.

\begin{example}
\label{example-global-sections}
Consider the rule which associates to every scheme $T$ the set
$F(T) = \Gamma(T, \mathcal{O}_T)$. We can turn this into a contravariant
functor by using for a morphism $f : T' \to T$ the pullback map
$f^\sharp : \Gamma(T, \mathcal{O}_T) \to \Gamma(T', \mathcal{O}_{T'})$.
Given a ring $R$ and an element $t \in R$ there exists a unique
ring homomorphism $\mathbf{Z}[x] \to R$ which maps $x$ to $t$.
Thus, using Lemma \ref{lemma-morphism-into-affine}, we see that
$$
\Mor(T, \Spec(\mathbf{Z}[x])) =
\Hom(\mathbf{Z}[x], \Gamma(T, \mathcal{O}_T)) =
\Gamma(T, \mathcal{O}_T).
$$
This does indeed
give an isomorphism $h_{\Spec(\mathbf{Z}[x])} \to F$.
What is the ``universal family'' $\xi$? To get it we have to
apply the identifications above to $\text{id}_{\Spec(\mathbf{Z}[x])}$.
Clearly under the identifications above this gives
that
$\xi = x \in \Gamma(\Spec(\mathbf{Z}[x]),
\mathcal{O}_{\Spec(\mathbf{Z}[x])}) = \mathbf{Z}[x]$
as expected.
\end{example}

\begin{definition}
\label{definition-representable-by-open-immersions}
Let $F$ be a contravariant functor on the category
of schemes with values in sets.
\begin{enumerate}
\item We say that $F$ {\it satisfies the sheaf property for the
Zariski topology} if for every scheme $T$ and every open covering
$T = \bigcup_{i \in I} U_i$, and for any collection of elements
$\xi_i \in F(U_i)$ such that $\xi_i|_{U_i \cap U_j} =
\xi_j|_{U_i \cap U_j}$ there exists a unique element
$\xi \in F(T)$ such that $\xi_i = \xi|_{U_i}$ in $F(U_i)$.
\item A {\it subfunctor $H \subset F$} is a rule that associates
to every scheme $T$ a subset $H(T) \subset F(T)$ such that
the maps $F(f) : F(T) \to F(T')$ maps $H(T)$ into
$H(T')$ for all morphisms of schemes $f : T' \to T$.
\item Let $H \subset F$ be a subfunctor. We say that
$H \subset F$ is {\it representable by open immersions}
if for all pairs $(T, \xi)$, where $T$ is a scheme and $\xi \in F(T)$
there exists an open subscheme $U_\xi \subset T$ with the following
property:
\begin{itemize}
\item[(*)] A morphism $f : T' \to T$ factors through $U_\xi$ if and only
if $f^*\xi \in H(T')$.
\end{itemize}
\item Let $I$ be a set. For each $i \in I$ let $H_i \subset F$
be a subfunctor. We say that the collection $(H_i)_{i \in I}$
{\it covers $F$} if and only if for every $\xi \in F(T)$
there exists an open covering $T = \bigcup U_i$ such that
$\xi|_{U_i} \in H_i(U_i)$.
\end{enumerate}
\end{definition}

\noindent
In condition (4), if $H_i \subset F$ is representable by open immersions
for all $i$, then to check $(H_i)_{i \in I}$ covers $F$, it suffices to
check $F(T) = \bigcup H_i(T)$ whenever $T$ is the spectrum of a field.

\begin{lemma}
\label{lemma-glue-functors}
Let $F$ be a contravariant functor on the category of schemes
with values in the category of sets. Suppose that
\begin{enumerate}
\item $F$ satisfies the sheaf property for the Zariski topology,
\item there exists a set $I$ and a collection of subfunctors
$F_i \subset F$ such that
\begin{enumerate}
\item each $F_i$ is representable,
\item each $F_i \subset F$ is representable by open immersions, and
\item the collection $(F_i)_{i \in I}$ covers $F$.
\end{enumerate}
\end{enumerate}
Then $F$ is representable.
\end{lemma}

\begin{proof}
Let $X_i$ be a scheme representing $F_i$ and let
$\xi_i \in F_i(X_i) \subset F(X_i)$ be the ``universal family''.
Because $F_j \subset F$ is representable by open immersions,
there exists an open $U_{ij} \subset X_i$ such that
$T \to X_i$ factors through $U_{ij}$ if and only if
$\xi_i|_T \in F_j(T)$. In particular
$\xi_i|_{U_{ij}} \in F_j(U_{ij})$ and therefore we obtain a
canonical morphism $\varphi_{ij} : U_{ij} \to X_j$ such that
$\varphi_{ij}^*\xi_j = \xi_i|_{U_{ij}}$. By definition of $U_{ji}$
this implies that $\varphi_{ij}$ factors through $U_{ji}$.
Since $(\varphi_{ij} \circ \varphi_{ji})^*\xi_j
=\varphi_{ji}^*(\varphi_{ij}^*\xi_j) =
\varphi_{ji}^*\xi_i = \xi_j$ we conclude that
$\varphi_{ij} \circ \varphi_{ji} = \text{id}_{U_{ji}}$
because the pair $(X_j, \xi_j)$ represents $F_j$.
In particular the maps $\varphi_{ij} : U_{ij} \to U_{ji}$
are isomorphisms of schemes.
Next we have to show that
$\varphi_{ij}^{-1}(U_{ji} \cap U_{jk}) = U_{ij} \cap U_{ik}$.
This is true because (a) $U_{ji} \cap U_{jk}$ is the largest
open of $U_{ji}$ such that $\xi_j$ restricts to an element
of $F_k$, (b) $U_{ij} \cap U_{ik}$ is the largest
open of $U_{ij}$ such that $\xi_i$ restricts to an element
of $F_k$, and (c) $\varphi_{ij}^*\xi_j = \xi_i$. Moreover,
the cocycle condition in Section \ref{section-glueing-schemes}
follows because both
$\varphi_{jk}|_{U_{ji} \cap U_{jk}} \circ
\varphi_{ij}|_{U_{ij} \cap U_{ik}}$ and
$\varphi_{ik}|_{U_{ij} \cap U_{ik}}$ pullback $\xi_k$
to the element $\xi_i$.
Thus we may apply Lemma \ref{lemma-glue-schemes}
to obtain a scheme $X$ with an open
covering $X = \bigcup U_i$ and isomorphisms $\varphi_i : X_i \to U_i$
with properties as in Lemma \ref{lemma-glue}.
Let $\xi_i' = (\varphi_i^{-1})^* \xi_i$.
The conditions of Lemma \ref{lemma-glue} imply that
$\xi_i'|_{U_i \cap U_j} = \xi_j'|_{U_i \cap U_j}$.
Therefore, by the condition that $F$ satisfies the sheaf
condition in the Zariski topology we see that
there exists an element $\xi' \in F(X)$ such that
$\xi_i = \varphi_i^*\xi'|_{U_i}$ for all $i$.
Since $\varphi_i$ is an isomorphism we also get that
$(U_i, \xi'|_{U_i})$ represents the functor $F_i$.

\medskip\noindent
We claim that the pair $(X, \xi')$ represents the functor $F$.
To show this, let $T$ be a scheme and let $\xi \in F(T)$.
We will construct a unique morphism $g : T \to X$ such that
$g^*\xi' = \xi$. Namely, by the condition that the subfunctors
$F_i$ cover $F$ there exists an open covering $T = \bigcup V_i$
such that for each $i$ the restriction $\xi|_{V_i} \in F_i(V_i)$.
Moreover, since each of the inclusions $F_i \subset F$ are representable
by open immersions we may assume that each $V_i \subset T$ is maximal
open with this property.
Because, $(U_i, \xi'|_{U_i})$ represents the functor $F_i$ we
get a unique morphism $g_i : V_i \to U_i$ such that
$g_i^*\xi'|_{U_i} = \xi|_{V_i}$. On the overlaps $V_i \cap V_j$
the morphisms $g_i$ and $g_j$ agree, for example because they both
pull back $\xi'|_{U_i \cap U_j} \in F_i(U_i \cap U_j)$
to the same element. Thus the morphisms $g_i$ glue to a unique morphism
from $T \to X$ as desired.
\end{proof}

\begin{remark}
\label{remark-representable-locally-ringed}
Suppose the functor $F$ is defined on all locally ringed spaces,
and if conditions of Lemma \ref{lemma-glue-functors} are replaced by
the following:
\begin{enumerate}
\item $F$ satisfies the sheaf property on the category of locally ringed
spaces,
\item there exists a set $I$ and a collection of subfunctors
$F_i \subset F$ such that
\begin{enumerate}
\item each $F_i$ is representable by a scheme,
\item each $F_i \subset F$ is representable by open immersions
on the category of locally ringed spaces, and
\item the collection $(F_i)_{i \in I}$ covers $F$
as a functor on the category of locally ringed spaces.
\end{enumerate}
\end{enumerate}
We leave it to the reader to spell this out further.
Then the end result is that the functor $F$ is
representable in the category of locally ringed spaces
and that the representing object is a scheme.
\end{remark}



